\documentclass[11pt,a4paper]{article}
\usepackage[margin=20mm]{geometry}
\usepackage{fontspec}
\setmainfont{DejaVu Serif}
\setsansfont{DejaVu Sans}
\setmonofont{DejaVu Sans Mono}[Scale=0.85]
\usepackage{polyglossia}
\setdefaultlanguage{russian}
\usepackage{graphicx}
\usepackage{hyperref}
\usepackage{enumitem}
\usepackage{fancyvrb}
\hypersetup{hidelinks}
\setlength{\parindent}{0pt}
\setlength{\parskip}{4pt}
\usepackage{titlesec}
\titleformat{\section}{\large\bfseries}{}{0pt}{}
\titlespacing*{\section}{0pt}{10pt}{4pt}
\sloppy
\setlist{nosep,leftmargin=*}
\emergencystretch=3em
\newcommand{\shot}[2]{\begin{center}\includegraphics[width=\linewidth,height=.34\textheight,keepaspectratio]{#1}\par\small #2\end{center}}
\begin{document}
\begin{center}
{\small Университет ИТМО · Фронт-энд разработка · 2026}\par
{\LARGE Домашняя работа № 3}\par
{\large CSS-переменные и темы}\par
Светлов Ярослав, группа J3211
\end{center}
\section*{Цель}
Добавить светлую и тёмную темы через CSS-переменные. За основу взята отдельная копия ДЗ2; предыдущие работы не заменяются.
\section*{Реализация}
\begin{itemize}
\item Переменные задают фон, текст, ссылки, контрастный фокус, поверхность логов и цвета графика.
\item theme.js отслеживает prefers-color-scheme. Тема меняется вслед за настройкой системы без перезагрузки страницы.
\item Атрибут data-bs-theme переключает цветовые переменные Bootstrap, чтобы формы и таблицы соответствовали общей теме.
\item Цвета Chart.js обновляются при смене темы. Отдельная кнопка не добавлена: выбран предусмотренный заданием вариант с системными настройками.
\end{itemize}

\section*{Проверка}
В Firefox 148 проверены кабинет, поиск, эксперимент и модели в обеих темах, смена предпочтения без перезагрузки, пустой результат и сброс фильтра, сохранение сессии. Результаты находятся в \texttt{checks/browser.json}.

В Chromium проверены вход, кабинет и график; ошибок консоли не обнаружено. В узком окне Firefox формы перестраиваются, таблицы прокручиваются внутри блока, страница не переполняется. Девять API-тестов пройдены.
\section*{Запуск и проверка}
Требуется Node.js 22.12 или новее. Команды выполняются из папки этой работы.
\begin{verbatim}
npm ci
npm start
\end{verbatim}
Открыть \url{http://localhost:3003}. Автотесты: \texttt{npm test}.
Учебные аккаунты: \texttt{admin@nexus.ai / password} и \texttt{vision@nexus.ai / cv12345}. Данные сохраняются в \texttt{db.local.json}, исходная база \texttt{db.json} не меняется.\section*{Как увидеть обе темы}
Открыть приложение и изменить светлую/тёмную тему в настройках системы либо эмулировать prefers-color-scheme в инструментах разработчика браузера. Перезагрузка страницы не требуется.
\section*{Вывод}
Обе цветовые схемы используют одну разметку. Цвета сосредоточены в переменных; логика запросов и сохранение данных не изменились.

\clearpage
\section*{Скриншоты}
\shot{checks/light.png}{Рисунок 1. Светлая тема.}
\shot{checks/dark.png}{Рисунок 2. Тёмная тема.}

\end{document}
